\begin{frame}{Suffix array}
  Dada una string de largo \(n\), imprimir los índices donde empiezan sus sufijos, en orden lexicográfico.

  \vspace{0.6em}
  \begin{columns}[c,onlytextwidth]
    \begin{column}{0.46\textwidth}
      \centering
      \renewcommand{\arraystretch}{1.2}
      \begin{tabular}{@{}c@{\quad}l@{}}
        índice & sufijo \\
        \(0\) & \texttt{caba} \\
        \(1\) & \texttt{aba} \\
        \(2\) & \texttt{ba} \\
        \(3\) & \texttt{a}
      \end{tabular}
    \end{column}
    \begin{column}{0.5\textwidth}
      Orden:\\[0.3em]
      \texttt{a}, \texttt{aba}, \texttt{ba}, \texttt{caba}\\[0.7em]
      Resultado: \(3\ 1\ 2\ 0\)
    \end{column}
  \end{columns}

  \vspace{0.8em}
  {\footnotesize \url{https://judge.yosupo.jp/problem/suffixarray}}
\end{frame}

\begin{frame}{Suffix array}
  Se arman ids para los substrings de largo 1.
  Esos se concatenan y se comprimen: salen los ids de largo 2, después los de largo 4, y así.

  \medskip
  Antes se agrega padding al final, para que los pares no se salgan de la string.

  \vspace{0.35em}
  \begin{center}
  \begin{tikzpicture}[x=11mm, y=7.4mm]
    \node[anchor=east, font=\small] at (-0.75,0) {letra};
    \node[anchor=east, font=\small] at (-0.75,-1.05) {id};
    \foreach \i/\ch in {0/c, 1/a, 2/b, 3/a} {
      \node[letter] at (\i,0) {\ch};
    }
    \foreach \i in {4,5,6} {
      \node[filler] at (\i,0) {\(\varepsilon\)};
    }
    \foreach \i/\id in {0/3, 1/1, 2/2, 3/1, 4/0, 5/0, 6/0} {
      \node[font=\small] at (\i,-1.05) {\id};
    }

    \node[anchor=east, font=\small] at (-0.75,-2.25) {par};
    \node[anchor=east, font=\small] at (-0.75,-3.35) {id};
    \foreach \i/\txt in {0/ca, 1/ab, 2/ba} {
      \node[letter, minimum width=9.6mm] at (\i+0.5,-2.25) {\txt};
    }
    \foreach \i/\txt in {3/{a\(\varepsilon\)}, 4/{\(\varepsilon\varepsilon\)}, 5/{\(\varepsilon\varepsilon\)}} {
      \node[filler, minimum width=9.6mm] at (\i+0.5,-2.25) {\txt};
    }
    \foreach \i/\id in {0/4, 1/2, 2/3, 3/1, 4/0, 5/0} {
      \node[font=\small] at (\i+0.5,-3.35) {\id};
    }
  \end{tikzpicture}
  \end{center}
\end{frame}

\begin{frame}{Suffix array}
  \(N\) es el largo de la string con el padding.

  \medskip
  \begin{center}
  \renewcommand{\arraystretch}{1.22}
  {\ttfamily
  \begin{tabular}{@{}l@{}}
    ids \(\leftarrow\) comprimir los caracteres \\
    \(l \leftarrow 1\) \\
    mientras \(l < n\): \\
    \quad par \(\leftarrow [(ids[i], ids[i+l])\) for \(i = 0\,..\!\!<\,N-2l+1]\) \\
    \quad ids \(\leftarrow\) comprimir par \\
    \quad \(l \leftarrow 2l\)
  \end{tabular}}
  \end{center}

  \vspace{0.55em}
  Al terminar, el sufijo que empieza en \(i\) queda en la posición \(\mathrm{ids}[i]\).
  El arreglo que hay que imprimir es la inversa: \(\mathrm{SA}[\mathrm{ids}[i]] = i\).

  \medskip
  En \texttt{caba}, después de la vuelta \(l = 2\), los ids son \(3,1,2,0\).
  La inversa es \(3\ 1\ 2\ 0\).
\end{frame}
